% Copyright (c) 2026 Eugene Joseph M. Ragasa.
% SPDX-License-Identifier: Apache-2.0
% Author: Eugene Joseph M. Ragasa
% Formal proof completed: 2026-10-27
\documentclass[11pt]{article}
\usepackage[T1]{fontenc}\usepackage{amsmath,amssymb}\usepackage[margin=1in]{geometry}\usepackage[hidelinks]{hyperref}
\title{Proof Dossier: Interval Upper-Bound Acceptance}\author{Eugene Joseph M. Ragasa}\date{Formal proof completed: 2026-10-27}
\begin{document}\maketitle
\section{Identity and status}
ID \texttt{BG-CERT-THM-017}; status \texttt{proved}. Lean declaration:
\par\noindent{\tiny\path{ProjectKoios.BandGauge.Certificates.}}\\
{\tiny\path{classifyUpperBound_eq_accepted_iff}}.
\section{Statement and proof}
An interval is accepted against an inclusive upper budget exactly when its endpoints are
ordered and its upper endpoint does not exceed the budget. The proof exhausts the three
ordered endpoint cases in the executable classifier.
\section{Boundary}
The result is source-neutral and assumes exact ordered inputs. It does not justify an
external enclosure, parser, or floating-point interpretation. Kernel axiom:
\texttt{propext}; no admissions. Checked with Lean \texttt{4.35.0} and Mathlib revision
\texttt{56a005855ea17e0b0958f73ec801ae140b0a1679}. AI-assisted and kernel-checked;
no independent expert review recorded.
\end{document}
